% Generated draft: requires mathematical and source review.
\documentclass[11pt]{article}
\usepackage[margin=25mm]{geometry}
\usepackage{fontspec}
\setmainfont{Latin Modern Roman}
\usepackage{amsmath,amssymb,mathtools,unicode-math}
\setmathfont{Latin Modern Math}
\usepackage{enumitem,needspace,xurl,hyperref}
\setlength{\parindent}{0pt}
\setlength{\parskip}{5pt}
\setlength{\emergencystretch}{3em}
\newcommand{\N}{\mathbb{N}}
\newcommand{\Z}{\mathbb{Z}}
\newcommand{\Q}{\mathbb{Q}}
\newcommand{\R}{\mathbb{R}}
\newcommand{\C}{\mathbb{C}}
\newcommand{\F}{\mathbb{F}}
\newcommand{\A}{\mathbb{A}}
\newcommand{\PP}{\mathbb{P}}
\newcommand{\E}{\mathbb{E}}
\newcommand{\eps}{\varepsilon}
\newcommand{\dd}{\,\mathrm{d}}
\begin{document}
\section*{UnsolvedMath ID: 30006579}
\textbf{Source problem code:} catalogue entry
\subsection{Definitions and mathematical statement}
Let \(p\) be a prime. Let \(K\) be either \(\Q\), equipped with the \(p\)-adic valuation, or \(k(\pi)\), equipped with the \(\pi\)-adic valuation, where \(k\) is a field of characteristic zero and \(\pi\) is an indeterminate. Write \(\widehat K\) for the completion of this valued field.

The supplied statement refers to an infinite-dimensional Tate ball \(B_\infty\) over \(\widehat K\) and a boundary-vanishing complex
\[
\mathcal C_\bullet=\widetilde{\Omega}_{\infty-\bullet}^{K\text{-}\mathrm{alg}}(B_\infty).
\]
Its coefficients are described as power series algebraic over \(K(z_1,z_2,\ldots)\). Here algebraic means satisfying a nonzero polynomial in one indeterminate with coefficients in that rational function field. The supplied record does not specify the power-series algebra, its convergence conditions, the construction of \(B_\infty\), the graded spaces and differential of \(\mathcal C_\bullet\), or the meaning of boundary vanishing. These remain source-dependent definitions, so the following is a conditional formulation using the intended complex.

\textbf{Conjecture 5 (as identified in the supplied source).} For every valued field \(K\) of either form above, the intended complex satisfies
\[
H_j(\mathcal C_\bullet)=0\qquad\text{for every nonzero degree }j,
\]
and its formal integration map
\[
I:H_0(\mathcal C_\bullet)\longrightarrow\widehat K
\]
is injective. Equivalently, for every \(\alpha\in H_0(\mathcal C_\bullet)\), one has \(I(\alpha)=0\Rightarrow\alpha=0\).

The damaged original excerpt associates this map with integration of a representative \(f\,\mathrm dz\) over \([0,1]^\infty\). This is formal integration notation; the supplied record does not define its infinite-dimensional meaning, convergence, or descent to homology. No analytic interpretation of that expression is imposed here. A fully self-contained formulation requires the missing definitions from the report.
\subsection{Short English statement}
For the stated \(p\)-adic and \(\pi\)-adic fields, the conjecture asks whether the intended complex of algebraic power-series forms vanishing on the boundary has zero homology outside degree zero, and whether formal integration detects every degree-zero homology class. The supplied record omits definitions needed to specify the complex and integration map completely.
\subsection{Sources}
\begin{itemize}
\item Supplied local UnsolvedMath registry record: permanent ID \(30006579\), title \emph{Adic Effective-Period Complex Conjecture}, exact source problem code \texttt{catalogue entry}; collection: Oberwolfach Reports Open Problems. The record supplies the year 2026 but names no individual proposer.
\item Supplied report citation: \emph{Model Theory and Arithmetic} (2026), Oberwolfach Reports, DOI 10.4171/owr/2026/3, pp.\ 157--216. \url{https://doi.org/10.4171/owr/2026/3}. The original excerpt labels the assertion Conjecture 5 and introduces it by analogy with Conjecture 2.
\item Supplied literature source: \emph{Oberwolfach Report 3/2026}, \url{https://doi.org/10.4171/owr/2026/3}. Literature assessment date: \texttt{2026-08-22}. This is supplied dated curation, not independently verified here. It labels the problem open, describes it as the adic analogue of a preceding period conjecture, and says that the report states the conjecture without proving the asserted acyclicity and integration-map property.
\item The registry labels the statement \texttt{corrected\_verified} and reports that the official report was checked directly while repairing OCR and notation. This is supplied verification metadata; no independent source inspection or mathematical verification is asserted in this draft.
\end{itemize}
\end{document}
